\section*{NS-86: optimized arithmetic potentials and a convolution correction}

\paragraph{Scope and prior art.}
This is an exact optimality reformulation, a complete tail estimate, and a
finite test of one nonlinear coefficient rule. No cofinal error contraction
is obtained. The integral-minus-lattice-sum operator is the classical
M\"untz operator, with the opposite sign and a joint cutoff required by
our logarithmic tail; see B\'aez-Duarte, arXiv:math/0505453, Sections 1--2.
The underlying approximation criterion and the zero-evaluation obstruction
are prior art. We do not claim that renaming either supplies new arithmetic
cancellation. The manuscript version remains v1.66.

\subsection*{1. Complete optimality in potential coordinates}
Keep the fixed $q=2$ space and the notation of NS-78:
\[
 R=\tau-\sum_{j\leq N}c_j\sigma_j,\quad r(t)=R(1/t),\quad
 \eta=\sum_{j\leq N}c_j/j,\quad A(z)=\int_0^z\{v\}\,dv/v.
\]
Thus $\sigma_j(1/t)=A(t/j)$, $\tau(1/t)=\log_+t$, and
$\|R\|^2=\int_0^\infty|r(t)|^2dt/t^2$.
Set
\[
 P(t)=\int_t^\infty \log(u/t)r(u)\,du/u^2,
 \qquad J(t)=\int_t^\infty r(u)\,du/u^2.
\]
Both integrals converge, $tP'(t)=-J(t)$, and $(tP'(t))'=r(t)/t^2$.
In particular,
\[
 |P(t)|\leq\sqrt{2\|R\|^2/t},\qquad
 |P'(t)|\leq\|R\|/t^{3/2}.
\]
For the actual finite residual, $r(t)=-\eta t$ below 1 and
$r(t)=u_0\log t+v_0+O(D/t)$ above $N$, where
\[
 u_0=1-\tfrac12\sum c_j,\quad
 v_0=\tfrac12\sum c_j\log(j/(2\pi)),\quad D=\tfrac16\sum j|c_j|.
\]
Consequently $P(t)=O(1+|\log t|^2)$ near zero and
$P(t)=(u_0\log t+2u_0+v_0)/t+O(D/(4t^2))$ at infinity.
No separate integral of $P$ over the full half-line is assumed finite.

\paragraph{Exact identity.}
For an integer $n\geq1$, define the joint-cutoff discrepancy
\[
 Q_{n,K}=\frac1n\int_0^{Kn}P(t)\,dt-\sum_{k=1}^K P(kn).
\]
Then
\begin{equation}
 \langle R,\sigma_n\rangle=\lim_{K\to\infty}Q_{n,K}.
 \tag{86.1}
\end{equation}
Indeed, distributionally,
$(t(A(t/n))')'=1/n-\sum_{k\geq1}\delta_{kn}$.
Twice integrating by parts gives (86.1). At zero both boundary products
are $O(t(1+|\log t|^2))$; at infinity they are
$O((1+\log^2t)/t)$. All interior joins are accounted for by the displayed
Dirac masses. Thus no endpoint or lattice term is discarded.

For the minimizing residual $R_N$, the exact normal equations become
$\lim_K Q_{n,K}=0$ for $n\leq N$. Moreover
\begin{equation}
 P_N(1)=\langle R_N,\tau\rangle=E_N.
 \tag{86.2}
\end{equation}
The second equality uses orthogonality to the fitted finite span. For a
nonoptimal residual it must not be substituted for $\|R\|^2$.

\paragraph{A complete observation-tail estimate.}
Write $T=Kn$ and $E_{>T}=\int_T^\infty|r(u)|^2du/u^2$. Cellwise
integration followed by absolutely convergent Fubini gives
\begin{equation}
 \langle R,\sigma_n\rangle-Q_{n,K}
 =-\int_T^\infty\{t/n\}P'(t)\,dt
 =\int_T^\infty\frac{r(u)}{u^2}
                 \bigl(A(u/n)-A(K)\bigr)\,du.
 \tag{86.3}
\end{equation}
Since $0\leq A(u/n)-A(K)\leq\log(u/T)$, Cauchy--Schwarz implies
\[
 |\langle R,\sigma_n\rangle-Q_{n,K}|
 \leq\sqrt{2E_{>T}/T}.
\]
For $K\geq1$, the complete NS-78 bound
$|A(v)-\tfrac12\log(2\pi v)|\leq1/(6v)$ improves this to
\begin{equation}
 |\langle R,\sigma_n\rangle-Q_{n,K}|
 \leq\sqrt{E_{>T}/T}\left(\frac1{\sqrt2}+\frac{n}{3T}\right).
 \tag{86.4}
\end{equation}
To check the constant, the logarithmic main term has squared norm
$1/(2T)$ in $L^2((T,\infty),du/u^2)$; the difference of its two
Stirling remainders has magnitude at most $1/(3K)$, hence norm at most
$1/(3K\sqrt T)$. Apply Minkowski, then Cauchy--Schwarz. This estimate
holds for each finite residual, with no bound on its coefficient mass.
It bounds an observation's tail, not that observation from below.

\paragraph{Tail cancellation before interval arithmetic.}
Let $P_T(t)=\int_t^T\log(u/t)r(u)du/u^2$. A useful exact finite formula is
\begin{equation}
 Q_{n,K}=\frac1n\int_0^T r(u)\,du/u
 -\sum_{k=1}^K P_T(kn)
 +\bigl[K(1-\log K)+\log K!\bigr]J(T).
 \tag{86.5}
\end{equation}
In its derivation the shared tail log-moment cancels identically; bounding
its two occurrences separately introduces unnecessary uncertainty.

For $T\geq N$, the remaining tail in (86.3) can also be evaluated with
the logarithmic model. Put $B=u_0\log n+2u_0+v_0$ and
$L_K=\sum_{k\leq K}(\log k)/k$. Its main value is
\[
 \frac1n\left[u_0\left(-\gamma_1-\tfrac12\log^2K+L_K\right)
       +B\left(H_K-\gamma-\log K\right)\right].
\]
Here $\gamma_1=\lim_K(L_K-\tfrac12\log^2K)$ is the first Stieltjes
constant. The error is at most $D/(4T^2)$: the derivative of the
potential remainder has magnitude at most $D/(2t^3)$, which is integrated
against $\{t/n\}\leq1$. This proof controls the entire tail.

\paragraph{Why the identity itself is not injectivity.}
For $1/2<\Re s<1$, the classical bounded Mellin evaluator in the physical
tail space is represented bilinearly by
\[
 r_s(t)=\begin{cases}t/(1-s),&t<1,\\t^{1-s},&t\geq1.\end{cases}
\]
Its potential is $t^{-s}/s^2$ above 1 and, with $L=\log(1/t)$,
$L^2/(2(1-s))+L/s+1/s^2$ below 1. The integral of this lower patch
equals the integral of $t^{-s}/s^2$ over $(0,1)$. Therefore
\[
 \lim_K Q_{n,K}=-\frac{\zeta(s)}{s^2n^s},\qquad P_s(1)=1/s^2.
\]
A hypothetical off-critical zero supplies a nonzero profile invisible to
all these observations. No such zero is asserted. This is the familiar
zero-evaluation obstruction in the new coordinates, not a new exclusion.
Also, NS-74's unitary multiplier commutes with dilations; covariance alone
does not distinguish its altered family from the original arithmetic one.

\subsection*{2. A specific arithmetic correction from $N$ to $N^2$}
Write $\mathbf1(n)=1$ and $e(n)=\mathbf1_{n=1}$ for arithmetic sequences,
and let $*$ denote Dirichlet convolution. Extend $c$ by zero after $N$.
The original divisor defect is $\delta=e+\mathbf1*c$. For $M=N^2$ define
\begin{equation}
 q_n=2c_n+\sum_{ij\mid n}c_i c_j\quad(1\leq n\leq M),\qquad
 q_n=0\quad(n>M),
 \tag{86.6}
\end{equation}
where $i,j\leq N$. At every $m\leq M$ one has exactly
\begin{equation}
 (e+\mathbf1*q)(m)=(\delta*\delta)(m)
                 =\sum_{d\mid m}\delta(d)\delta(m/d).
 \tag{86.7}
\end{equation}
Expand $(e+\mathbf1*c)^2$ in the associative convolution algebra.
Truncating $q$ cannot affect divisors of $m\leq M$, which proves the claim.
This is a convolution square, not a pointwise square, an energy square,
or a positivity statement. Beyond $M$ its omitted forcing is
\[
 (\delta*\delta)(m)-(e+\mathbf1*q)(m)
 =\sum_{\substack{d\mid m\\d>M}}(\mathbf1*c*c)(d).
\]
For the sharp canonical input $c_n=-\mu(n)$, $n\leq N$, (86.7) simply
reproduces the sharp canonical coefficients through $N^2$: each factorization
of $m\leq N^2$ has a factor at most $N$. It is not a new cure for the known
sharp-sequence obstruction, which does not itself rule out selected subsequences.

\paragraph{The three signed normalization sums are retained.}
Set $K_{ij}=\lfloor M/(ij)\rfloor$. Rearranging finite sums gives
\begin{align*}
 \eta_q&=2\eta_c+\sum_{i,j\leq N}\frac{c_ic_j}{ij}H_{K_{ij}},\\
 \sum_{n\leq M}q_n&=2\sum c_n+\sum_{i,j\leq N}c_ic_jK_{ij},\\
 \sum_{n\leq M}q_n\log n
 &=2\sum c_n\log n+
   \sum_{i,j\leq N}c_ic_j\bigl[K_{ij}\log(ij)+\log K_{ij}!\bigr].
\end{align*}
These give the exterior energy $\eta_q^2$ and the complete logarithmic
tail coefficients of the new residual. Orthogonality of the old residual
does not assert that these nonlinear sums are small.

\paragraph{Complete finite test, including old-space refitting.}
Let $V=\sum_{n\leq M}(q_n-c_n)\sigma_n$, $h=\langle R_N,V\rangle$,
$B=\|V\|^2$, and $B_\perp=\|(I-\Pi_N)V\|^2$. The unit-step error is
$E_N-2h+B$. The best real scalar along $V$ gains $h^2/B$.
Refitting all old coefficients and optimizing the scalar gains
$h^2/B_\perp$. These use the true complete Gram, not the elementary model.
After damping or old-space refitting, the modified coefficient vector
generally no longer satisfies the exact convolution-square identity (86.7).

\begin{center}
\begin{tabular}{rrrrr}
$N$&$M$&unit error$/E_N$&best raw gain/full gain&best projected gain/full gain\\
4&16&39.60315&0.16650&0.84935\\
8&64&0.47795&0.84502&0.84937\\
16&256&29.54956&0.10134&0.27307
\end{tabular}
\end{center}
The decimals are rounded outward-enclosed finite results, not predictions
for larger sizes. The N=16 unit step raises the error from about
$3.19477\cdot10^{-5}$ to $9.44040\cdot10^{-4}$. Its new energy beyond
$t=256$ alone exceeds 21 times the entire old error; its interior energy
also increases. Thus the failure is not solely an exterior-normalization
defect. A damped, old-space-refitted direction gives descent but captures
only 27.31 percent of available gain in this case.

The same coefficient rule is tested on the NS-74 altered loads with
$\beta=1/2+10^{-6}$. Its coefficient convolution identity remains algebraically
true, but (86.7) is not the physical divisor identity of the altered atoms.
That family still has its proved positive error floor. It must not be used
as evidence of a floor for the original family.

\subsection*{3. Outcome and stop condition}
The exact potential formulas and their complete tail bound are retained.
The proposed nonlinear correction fails the finite unit-step test at two
of three sizes; neither its eventual behavior nor all arithmetic correction
rules are ruled out. Optimizing a scalar does not supply a cofinal gain.
We have no new bound proving $E_{N^2}\leq\theta E_N$ for all sufficiently
large $N$ and fixed $\theta<1$. Such a bound would imply convergence;
NS-83 requires $\theta\geq1/2$, and still slower convergence could suffice.
The missing step remains a lower estimate on actual signed arithmetic
observations with their true cost. RH, G2, and the uniform Weil floor stay open.
